Bhagwat, N., Barry, A., Dickie, E. W., Brown, S. T., Devenyi, G. A., Hatano, K., ... \& Poline, J. B. (2021). Understanding the impact of preprocessing pipelines on neuroimaging cortical surface analyses. GigaScience, 10(1), giaa155.

\rev{Billot, B., Greve, D. N., Puonti, O., Thielscher, A., Van Leemput, K., Fischl, B., Dalca, A. V., \& Iglesias, J. E. (2023). SynthSeg: Segmentation of brain MRI scans of any contrast and resolution without retraining. Medical Image Analysis, 86, 102789.}

Botvinik-Nezer, R., Holzmeister, F., Camerer, C. F., Dreber, A., Huber, J., Johannesson, M., ... \& Rieck, J. R. (2020). Variability in the analysis of a single neuroimaging dataset by many teams. Nature, 582(7810), 84-88.

Bowring, A., Maumet, C., \& Nichols, T. E. (2019). Exploring the impact of analysis software on task fMRI results. Human brain mapping, 40(11), 3362-3384.

Chatelain, Y., Tetrel, L., Markiewicz, C. J., Goncalves, M., Kiar, G., Esteban, O., ... \& Glatard, T. (2023). A numerical variability approach to results stability tests and its application to neuroimaging. arXiv preprint arXiv:2307.01373.

\rev{Chatelain, Y., Sokolowski, A., Sharp, M., Poline, J.-B., \& Glatard, T. (2026). The practical impact of numerical variability on structural MRI measures of Parkinson's disease. Scientific Reports. https://doi.org/10.1038/s41598-026-64679-2}

Chougar, L., Faouzi, J., Pyatigorskaya, N., Yahia‐Cherif, L., Gaurav, R., Biondetti, E., ... \& Lehéricy, S. (2021). Automated categorization of parkinsonian syndromes using magnetic resonance imaging in a clinical setting. Movement Disorders, 36(2), 460-470.

Filip, P., Bednarik, P., Eberly, L. E., Moheet, A., Svatkova, A., Grohn, H., ... \& Mangia, S. (2022). Different FreeSurfer versions might generate different statistical outcomes in case–control comparison studies. Neuroradiology, 1-9.

Gaurav, R., Yahia‐Cherif, L., Pyatigorskaya, N., Mangone, G., Biondetti, E., Valabrègue, R., ... \& Lehéricy, S. (2021). Longitudinal changes in neuromelanin MRI signal in Parkinson's disease: a progression marker. Movement Disorders, 36(7), 1592-1602.

Ghosh, S. S., Poline, J. B., Keator, D. B., Halchenko, Y. O., Thomas, A. G., Kessler, D. A., \& Kennedy, D. N. (2017). A very simple, re-executable neuroimaging publication. F1000Research, 6.

Gomez-Ramirez, J., Quilis-Sancho, J., \& Fernandez-Blazquez, M. A. (2022). A comparative analysis of mri automated segmentation of subcortical brain volumes in a large dataset of elderly subjects. Neuroinformatics, 20, 63-72.

Gronenschild, E. H., Habets, P., Jacobs, H. I., Mengelers, R., Rozendaal, N., Van Os, J., \& Marcelis, M. (2012). The effects of FreeSurfer version, workstation type, and Macintosh operating system version on anatomical volume and cortical thickness measurements. PloS one, 7(6), e38234.

Haddad, E., Pizzagalli, F., Zhu, A. H., Bhatt, R. R., Islam, T., Ba Gari, I., ... \& Jahanshad, N. (2023). Multisite test–retest reliability and compatibility of brain metrics derived from FreeSurfer versions 7.1, 6.0, and 5.3. Human Brain Mapping, 44(4), 1515-1532.

Hanganu, A., Bedetti, C., Degroot, C., Mejia-Constain, B., Lafontaine, A. L., Soland, V., ... \& Monchi, O. (2014). Mild cognitive impairment is linked with faster rate of cortical thinning in patients with Parkinson’s disease longitudinally. Brain, 137(4), 1120-1129.

\rev{Hoffmann, M., Billot, B., Greve, D. N., Iglesias, J. E., Fischl, B., \& Dalca, A. V. (2022). SynthMorph: Learning contrast-invariant registration without acquired images. IEEE Transactions on Medical Imaging, 41(3), 543-558.}

\rev{Hoopes, A., Mora, J. S., Dalca, A. V., Fischl, B., \& Hoffmann, M. (2022). SynthStrip: Skull-stripping for any brain image. NeuroImage, 260, 119474.}

\rev{Klein, A., \& Tourville, J. (2012). 101 labeled brain images and a consistent human cortical labeling protocol. Frontiers in Neuroscience, 6, 171.}

Le, H., Zeng, W., Zhang, H., Li, J., Wu, X., Xie, M., ... \& Shen, J. (2020). Mean apparent propagator MRI is better than conventional diffusion tensor imaging for the evaluation of Parkinson’s disease: a prospective pilot study. Frontiers in Aging Neuroscience, 12, 563595.

Mak, E., Bergsland, N., Dwyer, M. G., Zivadinov, R., \& Kandiah, N. J. A. A. J. N. (2014). Subcortical atrophy is associated with cognitive impairment in mild Parkinson disease: a combined investigation of volumetric changes, cortical thickness, and vertex-based shape analysis. American Journal of Neuroradiology, 35(12), 2257-2264.

Mak, E., Su, L., Williams, G. B., Firbank, M. J., Lawson, R. A., Yarnall, A. J., ... \& O’Brien, J. T. (2015). Baseline and longitudinal grey matter changes in newly diagnosed Parkinson’s disease: ICICLE-PD study. Brain, 138(10), 2974-2986.

Masouleh, S., Eickhoff, S. B., Zeighami, Y., Lewis, L. B., Dahnke, R., Gaser, C., ... \& Valk, S. L. (2020). Influence of processing pipeline on cortical thickness measurement. Cerebral Cortex, 30(9), 5014-5027.

Mitchell, T., Lehéricy, S., Chiu, S. Y., Strafella, A. P., Stoessl, A. J., \& Vaillancourt, D. E. (2021). Emerging neuroimaging biomarkers across disease stage in Parkinson disease: a review. JAMA neurology, 78(10), 1262-1272.

Niso, G., Botvinik-Nezer, R., Appelhoff, S., De La Vega, A., Esteban, O., Etzel, J. A., ... \& Rieger, J. W. (2022). Open and reproducible neuroimaging: from study inception to publication. NeuroImage, 119623.

Rane, S., Plassard, A., Landman, B. A., Claassen, D. O., Donahue, M. J., \& Alzheimer’s Disease Neuroimaging Initiative. (2017). Comparison of cortical and subcortical measurements in normal older adults across databases and software packages. Journal of Alzheimer's disease reports, 1(1), 59-70.

Reuter M, Fischl B. Avoiding asymmetry-induced bias in longitudinal image processing. Neuroimage. 2011;57(1):19–21.

Reuter M, Rosas HD, Fischl B. Highly accurate inverse consistent registration: A robust approach. Neuroimage. 2010;53(4):1181–96.

Reuter M, Schmansky NJ, Rosas HD, Fischl B. Within-subject template estimation for unbiased longitudinal image analysis. Neuroimage. 2012;61(4):1402–18.

Sanz-Robinson, J., Torabi, M., Wishard, T. J., Glatard, T., Poline, J. B. (2023). Is between-pipeline variability affected by image quality? [Conference presentation]. OBHM 2023, Montreal, QC, Canada. www.humanbrainmapping.org

Sarasso, E., Agosta, F., Piramide, N., \& Filippi, M. (2021). Progression of grey and white matter brain damage in Parkinson's disease: a critical review of structural MRI literature. Journal of neurology, 268, 3144-3179.

Tessa, C., Lucetti, C., Giannelli, M., Diciotti, S., Poletti, M., Danti, S., ... \& Toschi, N. (2014). Progression of brain atrophy in the early stages of Parkinson's disease: A longitudinal tensor‐based morphometry study in de novo patients without cognitive impairment. Human brain mapping, 35(8), 3932-3944.

Wilson, H., Niccolini, F., Pellicano, C., \& Politis, M. (2019). Cortical thinning across Parkinson's disease stages and clinical correlates. Journal of the Neurological Sciences, 398, 31-38.
